
\subsection{RNS Base Conversion}
\label{sec:rns-base-conversion}
\begin{itemize}
    \item Introduce/justify why base conversion is important for the thesis
    \item Required during keyswitching etc.
    \item I start by introducing some terminology
\end{itemize}

\begin{definition}[RNS Basis.] 
    Given a large modulus $Q = \prod_{i=0}^{k-1} q_i$ composed of pairwise coprime factors, we define the \textit{RNS basis} $\mathcal{B}_Q = \{q_0, q_1, \dots, q_{k-1}\}$ as the set of RNS moduli used in the RNS/CRT representation \eqref{form:crt}.
\end{definition}

%%
%% BASE EXTENSION (NAIVE)
\iftrue{ % AI Written: Claude Fable
    % \color{purple}
    \begin{definition}[Base conversion]
        Let $\mathcal{B}_Q = \{q_0,\dots,q_{k-1}\}$ and
        $\mathcal{B} = \{b_0,\dots,b_{\ell-1}\}$ be RNS bases with
        $\gcd(Q, B) = 1$, where $B = \prod_{j} b_j$. Base conversion is
        the map that takes the representation of $x$ in $\mathcal{B}_Q$
        to the representation of $[x]_Q$ in $\mathcal{B}$, i.e.
        \[
            \big([x]_{q_0}, \dots, [x]_{q_{k-1}}\big)
            \;\longmapsto\;
            \big([x]_{b_0}, \dots, [x]_{b_{\ell-1}}\big).
        \]
        When the residues in $\mathcal{B}$ are appended to those in $\mathcal{B}_Q$, so that $x$ becomes represented in the \textit{extended} basis $\mathcal{B}_Q \cup \mathcal{B}$ for the modulus $QB$, the operation is also referred to as \emph{base extension}.
    \end{definition}
}\fi

% \begin{definition}[Base Extension.]
%     Special case of base conversion where we convert from $\mathcal{B}_Q \longmapsto \mathcal{B}_{QP}$. Note $\mathcal{B}_{QP} = \mathcal{B}_{Q} \cup \mathcal{B}_{P}$ and $P = \prod_i p_i$ etc. 
% \end{definition}

\begin{remark}
    Extend mod notation for RNS polynomials: $[a]_{\mathcal{B}_Q} = ([a]_{q_0}, \dots, [a]_{q_{k-1}})$ for the basis $\mathcal{B}_Q = \{q_i\}_{i=0}^{k-1}$ means $a$ must be in RNS form (\textcolor{red}{TODO: distinguish CRT vs NTT mode... maybe $[a]^\mathsf{NTT}_{\mathcal{B}_Q}$?})
\end{remark}

\begin{itemize}
    \item The naive approach to base conversion/extension: reconstruct $a$ and decompose it into the new basis
    \item Approximation exists called \textit{fast base extension}
\end{itemize}

%%%
%%% Base Extension
%%%
% \begin{algorithm}[ht!]
% \caption{Base Extension}
% \label{alg:base_extension}
% \begin{algorithmic}[1]
%     \Require $[a]_{\mathcal{B}_Q}^\mathsf{NTT}$ and a basis $\mathcal{B}_P$
%     \Ensure $[a]_{\mathcal{B}_{QP}}^\mathsf{NTT}$
%     \State Initialize an empty array $a_{P}$ of size $|\mathcal{B}_P|$
%     \Comment{Holds the $P$-towers}
%     \State $a_{\mathsf{coeffs}} := \mathsf{Inverse CRT}(\mathsf{Inverse NTT}(a))$
%     \Comment{Convert to coefficient form}
%     \For{$p_i \in \mathcal{B}_P$}
%         \State $a_{p_i} := a_{\mathsf{coeffs}} \pmod{p_i}$ \Comment{Find $p_i$-tower}
%         \State Append $\mathsf{NTT}(a_{p_i})$ to $a_{P}$
%     \EndFor
%     \State\Return $a \cup a_{P}$
% \end{algorithmic}
% \end{algorithm}

% \begin{enumerate}
%     \item Convert to RNS-coefficient mode, then convert to canonical/coefficient form
%     \item Find the residues/towers $\bmod{p_i}$ for $\mathcal{B}_P$
%     \item Apply NTT to each residue ($\bmod{p_i}$)
%     \item Append the NTT'd $P$-towers to the original $Q$-towers
% \end{enumerate}

\subsubsection{Fast Base Conversion}
\begin{itemize}
    \item Approximation to base extension introduced by \cite{cryptoeprint:2016/510}
    \item Performs the base extension in CRT mode, avoiding the CRT reconstruction + decomposition and multi-precision integers
    \item Produces not $x$ but $x + \alpha_x Q$ for some integer $\alpha_x \in [0, k)$
    % \item This error can be corrected..?
\end{itemize}

\begin{equation}
    \label{eq:fastbconv}
    \mathsf{FastBConv}(x, \mathcal{B}_Q, \mathcal{B}) = \left(\left[\sum_{i=0}^{k-1}\left[x_i\left(\frac{Q}{q_i}\right)^{-1}\right]_{q_i} \times \frac{Q}{q_i}\right]_b\right)_{b\in \mathcal{B}}
\end{equation}


\subsubsection{Fast Base Extension}
\begin{itemize}
    \item We can use \eqref{eq:fastbconv} as a subroutine in base extension
    \item Calculate only $\mathcal{B}_Q \longmapsto \mathcal{B}_P$, then use the fact that $\mathcal{B}_{QP} = \mathcal{B}_Q\cup \mathcal{B}_P$
\end{itemize}

%%%
%%% Fast Base Extension
%%%
\begin{algorithm}[ht!]
\caption{Fast Base Extension}
\label{alg:fast_base_extension}
\begin{algorithmic}[1]
    \Require $[a]_{\mathcal{B}_Q}$ and a basis $\mathcal{B}_P$
    \Ensure $[a + \epsilon]_{\mathcal{B}_{QP}}$, where $\epsilon = \alpha_aQ$ for some $\alpha_a\in[0, k)$
    \State $a_{Q} := \mathsf{Inverse NTT}(a)$
    \Comment{Convert to CRT form}
    \State $a_{P} := \mathsf{FastBConv(a_Q, \mathcal{B}_Q, \mathcal{B}_P)}$
    \State\Return $a \cup \mathsf{NTT}(a_{P})$
\end{algorithmic}
\end{algorithm}

\begin{itemize}
    \item Extension advantage over conversion:
    \begin{itemize}
        \item \eqref{eq:fastbconv} is approximate as a \textit{conversion} but exact as an \textit{extension}: the extended residue vector over $\mathcal{B}_Q\cup\mathcal{B}_P$ is an exact RNS representation of the lifted integer $x' = [x]_Q + \alpha_x Q \in \mathbb{Z}_{QP}$, since $\alpha_x Q \equiv 0 \pmod{q_i}$ keeps the original $Q$-residues consistent with $x'$
        \item The overflow $\alpha_x Q$ is removed (or shrunk to the small $\alpha_x$) by a subsequent division during mod down
        \item $\mathcal{B}_Q$ and $\mathcal{B}_P$ disjoint ($\gcd{(Q,P)} = 1)$ so the $Q$-residues can be reused
    \end{itemize}
\end{itemize}

% \iftrue{
%     \color{red}
%     3. The fast variant is approximate as a conversion but exact as an extension. This is the deepest justification, and it's worth a remark in your text. FastBConv doesn't compute [x]Q mod bj[x]_Q \bmod b_j
%     [x]Q​modbj​; it computes (x+uQ) mod bj(x + uQ) \bmod b_j
%     (x+uQ)modbj​ for some small overflow uu
%     u. Viewed as a conversion, that's an error. But viewed as an extension, the combined residue vector over BQ∪B\mathcal{B}_Q \cup \mathcal{B}
%     BQ​∪B is a perfectly consistent, exact RNS representation of the single lifted integer x′=[x]Q+uQ∈ZQBx' = [x]_Q + uQ \in \mathbb{Z}_{QB}
%     x′=[x]Q​+uQ∈ZQB​ — the original residues still agree with x′x'
%     x′ modulo each qiq_i
%     qi​, since uQuQ
%     uQ vanishes there. The RNS-variant schemes are designed around exactly this: they tolerate carrying x′x'
%     x′ instead of xx
%     x because a subsequent division by QQ
%     Q shrinks uQuQ
%     uQ to the small quantity uu
%     u, or because a Shenoy–Kumaresan-style correction (your FastBConvEx) removes it using one extra modulus. So "extension" isn't just the more common use case — it's the semantics under which the fast algorithm is actually correct.
% }\fi

\begin{remark}
    The factors $[(\frac{Q}{q_i})^{-1}]_{q_i}$ and $\frac{Q}{q_i}$ can be pre-calculated and stored in memory (will be discussed in \autoref{sec:impl:hps}).
\end{remark}


% This conversion is relatively fast. This is because the sum should ideally be
% reduced modulo q in order to provide the exact value x; instead, (2) provides
% x + αxq for some integer αx ∈ [0, k − 1]. Computing αx requires costly operations
% in RNS. So this step is by-passed, at the cost of an approximate result.
% FastBconv naturally extends to polynomials of R by applying it coefficient-
% wise.

% , which requires no such conversion but does result in some noise. \eqref{eq:fastbconv} converts $x$ from base $Q$ to base $B$. Using this we can convert the residues of $a \in R_Q$ to residues in $R_P$ and simply add the towers to \eqref{form:eval}.

\subsubsection{Modulus Switching}
\begin{itemize}
    \item Also referred to as \textit{mod down}
\end{itemize}

In RNS, switching from $Q$ to $Q' = Q/q_{k-1}$ amounts to \textit{dropping the last limb} with a small correction, computed limb-wise without CRT reconstruction:
\begin{equation}
    \label{eq:rns-modswitch}
    [c']_{q_i} = \left[q_{k-1}^{-1}\left(c - \delta\right)\right]_{q_i}, \quad 0 \leq i < k-1,
\end{equation}
where $\delta \equiv c \pmod{q_{k-1}}$ and, for BGV, $\delta \equiv 0 \pmod{t}$ so that decryption correctness is maintained. The noise shrinks by a factor $q_{k-1}$, matching the classical modulus switch.

\subsubsection{Approximate Mod Down $QP \to Q$}
\begin{itemize}
    \item Uses \eqref{eq:fastbconv} to convert the $P$-part back to $\mathcal{B}_Q$ before dividing by $P$:
    \begin{equation}
        \label{eq:approxmoddown}
        \mathsf{ModDown}(x) = \left[P^{-1}\left(x - \mathsf{FastBConv}([x]_{\mathcal{B}_P}, \mathcal{B}_P, \mathcal{B}_Q)\right)\right]_{\mathcal{B}_Q}
    \end{equation}
    \item For BGV the converted term must additionally be corrected modulo $t$ to preserve the plaintext \cite{cryptoeprint:2021/204}
    \item Corrects the fast base extension error! The conversion overflow $\alpha P$ is divided by $P$, shrinking it to the small additive term $\alpha$ --- this is why the approximation in \eqref{eq:fastbconv} is tolerable
    \item \autocite{cryptoeprint:2012/099, cryptoeprint:2018/117}
\end{itemize}
